\section{Вызовы авторизации}\label{ACME.Chall}

\subsection{Обработка вызова}\label{ACME.Chall.Proc}

С помощью вызовов авторизации сервер проверяет, что владелец учетной записи
контролирует определенный идентификатор. Сервер размещает вызовы в массиве
\sstr{challenges} объекта авторизации. Клиент отрабатывает вызов, демонстрируя
контроль над идентификатором. После отработки клиент дает ответ авторизации,
обращаясь по адресу, указанному в вызове. Сервер проверяет результаты
отработки.

В зависимости от типа идентификатора данная схема уточняется в трех аспектах:
содержание вызова авторизации, содержание отработки и содержание ответа
авторизации.

Вызов авторизации представляет собой объект JSON со следующими основными
полями:
\begin{itemize}
\item
\sstr{type} (обязательное, строка)~--- тип вызова;

\item
\sstr{url} (обязательное, строка)~--- URL-адрес, на который следует отправить 
ответ;

\item
\sstr{status} (обязательное, строка)~--- статус вызова. Возможные значения: 
\sstr{pending}, \sstr{processing}, \sstr{valid} и \sstr{invalid} 
(см.~\ref{ACME.Res.Status});

\item
\sstr{validated} (необязательное, строка)~--- момент времени, когда сервер 
проверил отработку вызова. Время задается в формате, указанном 
в~\cite{RFC3339}. Поле \sstr{validated} является обязательным, если поле 
\sstr{status} принимает значение \sstr{valid};

\item
\sstr{error} (необязательное, объект)~--- описание ошибки, которая возникла во
время проверки отработки вызова, если таковая имелась. Описание представляет 
собой объект с подробностями проблемы (см.~\ref{ACME.Transport.Err}). Несколько 
ошибок могут быть описаны с помощью подзадач. Вызов с ошибкой должен иметь 
статус \sstr{processing} или \sstr{invalid}.
\end{itemize}

% use: https://errata.rfc-editor.org/eid5732/

Дополнительные поля определяются типом вызова. 

Если сервер устанавливает в поле \sstr{status} значение \sstr{invalid}, 
то ему следует также включить в вызов поле \sstr{error}. 

Различные вызовы позволяют серверу проверять различные аспекты контроля над 
идентификатором. Выбор тех или иных вызовов определяется политикой сервера.
%
% skip: In some challenges, like HTTP
%  and DNS, the client directly proves its ability to do certain things
%  related to the identifier.  

В настоящем стандарте определяются правила проверки контроля над доменными
именами. Правила контроля над другими идентификаторами могут определяться 
в расширениях ACME. При расширении перечня типов идентификаторов должны 
определяться соответствующие вызовы авторизации.

% skip: 8.2 Retrying Challenges
% skip: https://errata.rfc-editor.org/eid7826/

\subsection{Вызов авторизации типа HTTP}\label{ACME.Chall.HTTP}

При использовании вызова типа HTTP клиент доказывает контроль над доменным 
именем, размещая ресурс HTTP по определенному сетевому адресу в области 
действия доменного имени.

% skip: As a domain may resolve to multiple IPv4 and IPv6 addresses, the
%   server will connect to at least one of the hosts found in the DNS A
%   and AAAA records, at its discretion.  Because many web servers
%   allocate a default HTTPS virtual host to a particular low-privilege
%   tenant user in a subtle and non-intuitive manner, the challenge must
%   be completed over HTTP, not HTTPS.

% skip: https://errata.rfc-editor.org/eid6843/

В поле \sstr{type} вызова авторизации типа HTTP указывается значение
\sstr{http-01}. \pdfexample{Вызов}{acme/chall1.txt} содержит следующее 
дополнительное поле:

\begin{itemize}
\item
\sstr{token} (обязательное, строка)~--- случайное значение, однозначно 
идентифицирующее вызов. Это значение должно содержать не менее 128 бит 
энтропии. Оно не должно содержать символы за пределами алфавита base64url и не 
должно включать завершающие символы \str{=}.
%
В частности, строка должна состоять из не менее чем 22 символов.
%
% note: log2(64^22) = 132
% use: https://errata.rfc-editor.org/eid6950/
\end{itemize}

\if0
\begin{codeblock}
{
  "type": "http-01",
  "url": "https://example.com/acme/chall/prV_B7yEyA4",
  "status": "pending",
  "token": "LoqXcYV8q5ONbJQxbmR7SCTNo3tiAXDfowyjxAjEuX0"
}
\end{codeblock}
\fi

Клиент \pdfexample{отрабатывает}{acme/chall2.txt} вызов, вычисляя ключ 
авторизации и размещая его по адресу, который получается конкатенацией 
фиксированного префикса \sstr{/.well-known/acme-challenge} со значением поля 
\sstr{token} вызова.

Ключ авторизации представляет собой строку, которая получается добавлением к 
значению поля \sstr{token} сначала точки (\str{.}), а затем отпечатка 
(thumbprint) открытого ключа учетной записи. Отпечаток вычисляется по 
правилам~\cite{RFC7638} c помощью алгоритма хэширования \code{belt-hash}. 
Хэш-значение представляется строкой в алфавите base64url.

% change: SHA-256 -> belt-hash
%
% skip: As noted in [RFC7518] any prepended zero octets in the fields 
% of a JWK object MUST be stripped before doing the computation.
%
% skip: https://errata.rfc-editor.org/eid7565/

\if0
\begin{codeblock}
GET /.well-known/acme-challenge/LoqXcYV8...jxAjEuX0 HTTP/1.1
Host: example.org

HTTP/1.1 200 OK
Content-Type: application/octet-stream

LoqXcYV8...jxAjEuX0.9jg46WB3...fm21mqTI
\end{codeblock}
\fi 

\pdfexample{Ответ}{acme/chall3.txt} авторизации клиента~--- это пустой объект 
JSON (\sstr|\{\}|), который сигнализирует, что вызов может быть проверен 
сервером.

\if0
\begin{codeblock}
POST /acme/chall/prV_B7yEyA4 HTTP/1.1
Host: example.com
Content-Type: application/jose+json

{
  "protected": base64url({
    "alg": "BIGNS128",
    "kid": "https://example.com/acme/acct/evOfKhNU60wg",
    "nonce": "UQI1PoRi5OuXzxuX7V7wL0",
    "url": "https://example.com/acme/chall/prV_B7yEyA4"
  }),
  "payload": base64url({}),
  "signature": "Q1bURgJoEslbD1c5...3pYdSMLio57mQNN4"
}
\end{codeblock}
\fi

Получив ответ, сервер независимо вычисляет ключ авторизации и 
проверяет контроль клиента над доменным именем следующим образом:

\begin{enumerate}
\item     
Сформировать URL-адрес, заполняя шаблон 
\sstr{http://\{domain\}/.well-known/acme-challenge/\{token\}}.
В поле \sstr{domain} сервер указывает проверяемое доменное имя,
в поле \sstr{token}~--- токен из вызова авторизации.
\item
Убедиться, что полученный URL-адрес имеет правильный формат.
\item
Разыменовать URL-адрес с помощью GET-запроса. Этот запрос должен быть 
отправлен на TCP-порт 80 HTTP-сервера.
% 
% skip: https://errata.rfc-editor.org/eid8381/
%   The HTTP client must not resolve and/or must ignore any HTTPS DNS RRs 
%   [RFC 9460].
%
\item
Убедиться, что тело ответа представляет собой ключ авторизации корректного 
формата. Сервер должен игнорировать пробелы в конце тела. 
\item
Убедиться, что предоставленные HTTP-сервером ключ авторизации соответствуют 
ключу, который вычислил сервер.
\end{enumerate}

% skip: The server SHOULD follow redirects when dereferencing the URL.
%  Clients might use redirects, for example, so that the response can be
%  provided by a centralized certificate management server.  See
%  Section 10.2 for security considerations related to redirects.

Отработка вызова завершена успешно тогда и только тогда, когда успешно 
завершены все перечисленные проверки.

Клиент должен отменить выделение ресурса для отработки вызова сразу после 
завершения отработки, т.~е. после того, как поле \sstr{status} вызова примет 
значение \sstr{valid} или \sstr{invalid}.

% skip: Note that because the token appears both in the request sent by the
%   ACME server and in the key authorization in the response, it is
%   possible to build clients that copy the token from request to
%   response.  Clients should avoid this behavior because it can lead to
%   cross-site scripting vulnerabilities; instead, clients should be
%   explicitly configured on a per-challenge basis.  A client that does
%   copy tokens from requests to responses MUST validate that the token
%   in the request matches the token syntax above (e.g., that it includes
%   only characters from the base64url alphabet).

\subsection{Вызов авторизации типа DNS}\label{ACME.Chall.DNS}

При использовании вызова типа DNS клиент доказывает контроль над доменным 
именем, размещая TXT-запись DNS на сервере, доступном по этому имени.

В поле \sstr{type} вызова авторизации типа HTTP указывается значение
\sstr{dns-01}. \pdfexample{Вызов}{acme/chall4.txt} содержит следующее 
дополнительное поле:

\begin{itemize}
\item
\sstr{token} (обязательное, строка)~--- случайное значение, 
однозначно идентифицирующее вызов. Это значение должно 
содержать не менее 128 бит энтропии. Оно не должно содержать символы за 
пределами алфавита base64url и не должно включать завершающие символы \str{=}. 
\end{itemize}

\if0
\begin{codeblock}
{
  "type": "dns-01",
  "url": "https://example.com/acme/chall/Rg5dV14Gh1Q",
  "status": "pending",
  "token": "evaGxfADs6pSRb2LAv9IZf17Dt3juxGJ-PCt92wr-oA"
}
\end{codeblock}
\fi

Клиент отрабатывает вызов, вычисляя ключ авторизации, хэшируя
его и размещая хэш-значение в TXT-записи DNS.
TXT-запись касается домена \code|_acme-challenge.{domain}|, где 
\code{domain}~--- проверяемое доменное имя.

Ключ авторизации строится так же, как в вызове типа HTTP.
Хэширование выполняется с помощью алгоритма \code{belt-hash}.
Хэш-значение кодируется по правилам base64url.

Пример записи DNS:

\begin{codeblock}
_acme-challenge.www.example.org. 300 IN TXT "gfj9Xq...Rg85nM"
\end{codeblock}

\if0
Например, если проверяемое доменное имя — "www.example.org", то 
клиент предоставит следующую DNS-запись : 

\begin{codeblock}
_acme-challenge.www.example.org. 300 IN TXT "gfj9Xq...Rg85nM"
\end{codeblock}
\fi

\pdfexample{Ответ}{acme/chall5.txt} авторизации клиента~--- это пустой объект 
JSON (\sstr|{}|), который сигнализирует, что вызов может быть проверен 
сервером.

\if0
\begin{codeblock}
POST /acme/chall/Rg5dV14Gh1Q HTTP/1.1
Host: example.com
Content-Type: application/jose+json

{
  "protected": base64url({
    "alg": "BIGNS128",
    "kid": "https://example.com/acme/acct/evOfKhNU60wg",
    "nonce": "SS2sSl1PtspvFZ08kNtzKd",
    "url": "https://example.com/acme/chall/Rg5dV14Gh1Q"
  }),
  "payload": base64url({}),
  "signature": "Q1bURgJoEslbD1c5...3pYdSMLio57mQNN4"
}
\end{codeblock} 
\fi

Получив ответ, сервер независимо вычисляет ключ авторизации и 
проверяет контроль клиента над доменным именем следующим образом:

\begin{enumerate}
\item
Вычислить хэш-значение ключа авторизации.
\item
Запросить TXT-записи DNS для проверяемого доменного имени.
\item
Убедиться, что в одной из TXT-записей указано вычисленное хэш-значение. 
\end{enumerate}

Отработка вызова завершена успешно тогда и только тогда, когда успешно 
завершены все перечисленные проверки.

Клиенту следует отменить выделение ресурса для отработки вызова сразу после 
завершения отработки, т.~е. после того, как поле \sstr{status} вызова примет 
значение \sstr{valid} или \sstr{invalid}.
